\documentclass[UTF8,11pt,a4paper]{ctexart}

\usepackage[margin=1in]{geometry}
\usepackage{amsmath,amssymb,amsthm,mathtools}
\usepackage{enumitem}
\usepackage{fancyhdr}
\usepackage{hyperref}
\usepackage{xcolor}

\hypersetup{
  colorlinks=true,
  linkcolor=blue!60!black,
  citecolor=blue!60!black,
  urlcolor=blue!60!black
}

\setlength{\parindent}{2em}
\setlength{\parskip}{0.4em}
\setlist[itemize]{leftmargin=2em}
\setlist[enumerate]{leftmargin=2em}
\setcounter{tocdepth}{1}

\newtheorem{definition}{定义}
\newtheorem{proposition}{命题}
\newtheorem{theorem}{定理}

\newcommand{\todo}[1]{\textcolor{red}{[TODO: #1]}}
\newcommand{\N}{\mathbb{N}}
\newcommand{\Prb}{\operatorname{Pr}}
\newcommand{\E}{\mathbb{E}}

\pagestyle{fancy}
\fancyhf{}
\fancyfoot[C]{\thepage}
\renewcommand{\headrulewidth}{0pt}
\renewcommand{\footrulewidth}{0pt}

\fancypagestyle{plain}{
  \fancyhf{}
  \fancyfoot[C]{\thepage}
  \renewcommand{\headrulewidth}{0pt}
  \renewcommand{\footrulewidth}{0pt}
}

\title{\Huge\textbf{Lower Bounds on Locality Sensitive Hashing Review}\\[0.5em]}
\author{241880622\quad 蒲今}
\date{\today}

\begin{document}

\maketitle

\vspace{-1.2em}
\begingroup
\small
\setlength{\parskip}{0pt}
\tableofcontents
\endgroup
\newpage

\section{背景与问题定义}

\subsection{近似最近邻问题}

近似最近邻问题要求：给定一个预处理后的点集 $P$，对于查询点 $q$，快速找到一个点
$x\in P$，使得
\[
  d(q,x)\le c\cdot \min_{y\in P} d(q,y),
\]
其中 $c>1$ 是近似因子。当 $c=1$ 时是精确最近邻；当 $c>1$ 时，算法允许返回一个足够接近最优答案的点。
在高维空间中，精确最近邻通常很难高效解决，因此近似最近邻成为更实际的研究对象。

\subsection{LSH 框架与复杂度指数}

实际算法通常使用近邻问题的判定版本：在已知数据库中存在点与查询点距离至多为 $r$ 的条件下，
返回一个距离至多为 $cr$ 的点。Indyk 和 Motwani 提出的局部敏感哈希
（Locality Sensitive Hashing, LSH）通过随机哈希使近点更容易碰撞、远点更难碰撞，
从而只检查少量候选点 \cite{im1998}。若近点碰撞概率至少为 $p$、远点碰撞概率至多为 $q$，则
\[
  \rho=\frac{\log(1/p)}{\log(1/q)}
\]
控制查询复杂度中的指数：查询通常需要约 $n^\rho$ 次距离计算，空间中则出现约 $n^{1+\rho}$ 的项。
因此，构造问题是尽量减小 $\rho$，而下界问题则要回答：在标准 LSH 模型内，$\rho$ 究竟还能小到什么程度？

\subsection{本文评述的问题}

本文评述 Motwani、Naor 和 Panigrahy 的
\emph{Lower Bounds on Locality Sensitive Hashing} \cite{mnp2006}。论文不提出新的哈希构造，
而是首次对整个标准 LSH 模型给出统一的非平凡下界，将“尚未找到更好的算法”推进为
“任何满足该模型定义的算法都受到共同限制”。

\section{主要定理与贡献}

\subsection{正式定义}

设 $(X,d_X)$ 是一个度量空间。随机哈希函数族 $H:X\to\N$ 称为
$(r,R,p,q)$-sensitive，如果对任意 $x,y\in X$，有
\[
  d_X(x,y)\le r
  \Rightarrow
  \Prb_H[H(x)=H(y)]\ge p,
\]
并且
\[
  d_X(x,y)>R
  \Rightarrow
  \Prb_H[H(x)=H(y)]\le q.
\]
在近似最近邻问题中通常令 $R=cr$，并用
\[
  \rho=\frac{\log(1/p)}{\log(1/q)}
\]
衡量 LSH 的效率。上式中的 $\rho$ 只衡量一个给定哈希族的质量。为了描述度量空间
$(X,d_X)$ 上所有合法哈希族能够达到的最佳指数，进一步定义
\[
  \rho_X(c)
  =
  \sup_{r>0}
  \inf\left\{
    \frac{\log(1/p)}{\log(1/q)}:
    \begin{array}{l}
      0<p,q<1,\ \text{且 }X\text{ 上存在一个}\\[-0.2em]
      (r,cr,p,q)\text{-sensitive 哈希族}
    \end{array}
  \right\}.
\]
对于高维 $\ell_s$ 空间，再定义其渐近最佳指数为
\[
  \rho_s(c)
  =
  \limsup_{d\to\infty}\rho_{\ell_s^d}(c).
\]
因此，$\rho$ 是单个哈希族的性能参数，而 $\rho_X(c)$ 和 $\rho_s(c)$ 分别表示在给定空间上
对所有合法哈希族优化后的最佳指数，以及该指数在维数趋于无穷时的渐近值。这里 $p$ 越大，
说明近点越容易被放到同一个桶中；$q$ 越小，说明远点越不容易误碰撞。因此，一个好的 LSH 需要同时做到
“近点高碰撞”和“远点低碰撞”。本文的技术正是围绕这两个要求之间的矛盾展开。

\subsection{主定理及适用范围}

论文的主定理指出，对任意 $s>0$ 和 $c>1$，
\[
  \rho_s(c)
  \ge
  \frac{e^{1/c^s}-1}{e^{1/c^s}+1}
  \ge
  \frac{0.462}{c^s}.
\]
特别地，$\ell_1$ 与 $\ell_2$ 空间分别得到 $\Omega(1/c)$ 和 $\Omega(1/c^2)$ 下界。
该结论约束的是数据无关的标准 LSH 哈希族，而不是所有近似最近邻数据结构；它也不直接排除
数据依赖哈希、图搜索或其他更强模型中的改进。论文的贡献在于首次从 LSH 定义本身推出统一下界，
并建立了一套后来可继续加强的 Hamming 立方体分析方法。

\section{证明总体思路}

\subsection{转化到 Hamming 立方体}

论文没有直接分析整个高维 $\ell_s^d$ 空间，而是先分析它的一个特殊子集：
\[
  \{0,1\}^d.
\]
这个集合是 Hamming 立方体。在该集合上，两个点之间的 $\ell_1$ 距离等于它们不同坐标的个数。

选择 Hamming 立方体的好处是，它既足够简单，又保留了高维空间的核心困难。立方体中每个点都是一个 0-1 向量，
距离就是不同坐标的数量；同时，从一个点出发随机翻转坐标，就得到一个非常自然的随机游走过程。论文正是利用这个随机游走来刻画
“从一个点移动到附近点”时哈希值是否仍然相同。

\subsection{随机桶与随机游走}

第一层是哈希桶分析。固定一个点 $x$，并从哈希族 $\mathcal{H}$ 中随机抽取 $H$，定义随机桶
\[
  A_H(x)=H^{-1}(H(x)).
\]
这个集合就是 $x$ 在随机哈希函数 $H$ 下所在的桶。远点低碰撞条件说明：如果一个点离 $x$ 很远，
那么对 $H$ 的随机性取概率，它落入 $A_H(x)$ 的概率至多是 $q$。因此，这一条件限制的是随机桶大小的期望，
而不是每一个确定哈希函数所产生的桶大小。

第二层是随机游走分析。论文考虑从点 $x$ 出发，在 Hamming 立方体上随机走 $r$ 步。由于每一步只翻转一个坐标，
终点与起点的 Hamming 距离不会超过 $r$。先对随机哈希函数使用近点碰撞条件，再对随机游走终点取平均，
可知终点以至少 $p$ 的联合概率落入随机桶 $A_H(x)$。

\subsection{小集合扩散与核心矛盾}

第三层是小集合扩散性质。Hamming 立方体上的随机游走会扩散：如果集合 $B$ 很小，那么从 $B$ 中出发的随机游走很难在若干步之后仍然停留在 $B$ 中。
这与 LSH 要求的近点高碰撞形成冲突。论文通过傅里叶分析给出这一随机游走概率的上界，并把它与哈希桶大小的上界结合，得到核心不等式
\[
  p
  \le
  \left(
    q+
    \exp\left(
      -\frac{1}{d}\left(\frac{d}{2}-R\right)^2
    \right)
  \right)^{
    \frac{e^{2r/d}-1}{e^{2r/d}+1}
  }.
\]
这个不等式说明，只要远点碰撞概率 $q$ 很小，近点碰撞概率 $p$ 也会被迫受到限制。
最后再把 Hamming 立方体看作 $\ell_s^d$ 的子集，就可以把该下界推广到高维 $\ell_s$ 空间。

证明的核心冲突可以概括为
\[
  \text{远点低碰撞要求桶不能太大}
  \quad+\quad
  \text{小桶难以困住随机游走}
  \quad\Rightarrow\quad
  \text{近点碰撞概率存在上限。}
\]

\section{关键技术细节}

\subsection{第一步：限制哈希桶大小}

固定一个点 $x\in\{0,1\}^d$，并从哈希族 $\mathcal{H}$ 中随机抽取哈希函数 $H$，定义随机桶
\[
  A_H(x)=H^{-1}(H(x)).
\]
集合 $A_H(x)$ 是与 $x$ 哈希到同一个值的所有点。注意，LSH 的碰撞概率是对 $H\sim\mathcal{H}$ 而言的；
因此，远点条件控制的是 $|A_H(x)|$ 对随机哈希函数的期望，而不是某个固定 $H$ 的桶大小。

为了估计桶的大小，可以把所有点 $u\in\{0,1\}^d$ 分成两类。第一类是距离 $x$ 不超过 $R$ 的点，
它们的数量最多是 Hamming 球的大小
\[
  \sum_{k=0}^{\lfloor R\rfloor}\binom{d}{k}.
\]
对这些点，碰撞概率最多只能粗略地用 $1$ 上界。第二类是距离 $x$ 大于 $R$ 的点。由 LSH 的远点条件，
它们与 $x$ 落入同一桶的概率至多是 $q$。对所有 $u$ 的碰撞指示变量求和，得到
\[
  \E_{H\sim\mathcal{H}}|A_H(x)|
  =
  \sum_{u\in\{0,1\}^d}
  \Prb_{H\sim\mathcal{H}}[H(u)=H(x)]
  \le
  \sum_{k=0}^{\lfloor R\rfloor}\binom{d}{k}
  +
  q\sum_{k=\lfloor R\rfloor+1}^{d}\binom{d}{k}.
\]
把它除以整个立方体大小 $2^d$，就得到桶的平均密度上界。进一步利用 Hamming 球体积估计，当 $R<d/2$ 时，
\[
  \E_{H\sim\mathcal{H}}\frac{|A_H(x)|}{2^d}
  \le
  q+
  \exp\left(
    -\frac{1}{d}\left(\frac d2-R\right)^2
  \right).
\]
这里第一项 $q$ 来自远点碰撞概率，第二项来自半径为 $R$ 的 Hamming 球本身在整个立方体中所占的比例。

\subsection{第二步：用随机游走连接近点碰撞}

在 Hamming 立方体上，从点 $x$ 出发随机翻转一个坐标，重复 $r$ 次，得到随机点 $W_r(x)$。
由于每一步只改变一个坐标，所以
\[
  \|x-W_r(x)\|_1\le r.
\]
因此，对随机游走产生的每一个终点 $y=W_r(x)$，近点条件都给出
\(\Prb_{H\sim\mathcal{H}}[H(y)=H(x)]\ge p\)。再对游走终点取平均，得
\[
  \Prb_{H\sim\mathcal{H},\,W_r}
  [H(W_r(x))=H(x)]
  \ge p.
\]
这里的概率同时包含两种相互独立的随机性：$H$ 从 $\mathcal{H}$ 中随机抽取，$W_r(x)$ 是随机游走的终点。
对任意固定的 $H$，碰撞事件等价于 $W_r(x)\in A_H(x)$；但 LSH 定义保证的是上述联合概率的下界，
并不保证对每个固定的 $H$，随机游走都以至少 $p$ 的概率留在桶中。

\subsection{第三步：用小集合扩散性质给出上界}

Hamming 立方体上的随机游走具有扩散性：如果集合 $B$ 不够大，随机游走从 $B$ 中出发后仍留在 $B$ 中的概率不会太高。
论文用 Fourier 分析和 Bonami-Beckner 不等式证明了如下随机游走引理：对任意非空集合 $B\subseteq\{0,1\}^d$，
如果 $Q_B$ 表示从 $B$ 中均匀随机选点并走 $r$ 步后的终点，则
\[
  \Prb[Q_B\in B]
  \le
  \left(\frac{|B|}{2^d}\right)^\alpha,
  \qquad
  \alpha=
  \frac{e^{2r/d}-1}{e^{2r/d}+1}.
\]
这个引理的含义是：桶的相对密度越小，随机游走留在桶中的概率越小。

这里 Fourier 分析的作用，是把“随机游走是否留在集合中”这个组合问题转化为对函数频率成分的估计。
具体地说，对集合 $B$ 考虑它的示性函数 $\mathbf{1}_B$。在 Hamming 立方体上，任何函数都可以用 Walsh 函数展开：
\[
  \mathbf{1}_B
  =
  \sum_{S\subseteq\{1,\dots,d\}}
  \widehat{\mathbf{1}_B}(S)W_S.
\]
这里 $S$ 可以理解为一组坐标，$W_S$ 描述这些坐标上的奇偶性，$\widehat{\mathbf{1}_B}(S)$ 则表示集合 $B$ 在这种“频率模式”上的权重。

随机游走矩阵 $P$ 对 Walsh 函数的作用非常简单：
\[
  P W_S
  =
  \left(1-\frac{2|S|}{d}\right)W_S.
\]
也就是说，Walsh 函数正好是随机游走的特征向量。经过 $r$ 步随机游走后，频率 $S$ 对应的权重会乘上
\[
  \left(1-\frac{2|S|}{d}\right)^r.
\]
不能笼统地说所有 $|S|$ 较大的成分都会被快速削弱，因为特征值
\[
  \lambda_S=1-\frac{2|S|}{d}
\]
在 $|S|$ 接近 $d$ 时接近 $-1$，其绝对值并不小。正确的处理需要分两部分。当 $|S|\le d/2$ 时，$\lambda_S\ge0$，且
\[
  0\le
  \left(1-\frac{2|S|}{d}\right)^r
  \le
  \exp\left(-\frac{2r|S|}{d}\right),
\]
因此这些非负特征值的贡献随 $|S|$ 增大而指数衰减。当 $|S|>d/2$ 时，$\lambda_S<0$；由于论文要求 $r$ 为奇数，
$\lambda_S^r$ 仍为负数，而其前面的 Fourier 系数平方非负，所以这些项在求上界时可以直接丢弃。
这正是该证明要求 $r$ 为奇数的实质原因。

于是，随机游走从 $B$ 出发后仍落在 $B$ 中的概率可以写成 Fourier 系数的加权和：
\[
  \Prb[Q_B\in B]
  =
  \frac{2^d}{|B|}
  \sum_{S\subseteq\{1,\dots,d\}}
  \widehat{\mathbf{1}_B}(S)^2
  \left(1-\frac{2|S|}{d}\right)^r.
\]
接下来需要控制这个加权和。Bonami-Beckner 不等式提供的正是这种控制：当 $B$ 的密度 $|B|/2^d$ 很小时，
它的 Fourier 系数不可能让上式过大。换句话说，小集合不能在随机游走下保持太强的稳定性。
这一步最终推出
\[
  \Prb[Q_B\in B]
  \le
  \left(\frac{|B|}{2^d}\right)^{
    \frac{e^{2r/d}-1}{e^{2r/d}+1}
  }.
\]

\subsection{第四步：用 Jensen 不等式闭合上下界}

令 $x$ 在 $\{0,1\}^d$ 上均匀分布，并与 $H$、随机游走相互独立。条件于固定的 $H$ 和
$H(x)=k$，起点 $x$ 在非空桶 $H^{-1}(k)$ 中均匀分布，因而可以对每个桶应用上述随机游走引理。记
\[
  \alpha=\frac{e^{2r/d}-1}{e^{2r/d}+1},
\]
则近点碰撞条件和随机游走引理共同给出
\[
  p
  \le
  \E_H\E_x
  \left(\frac{|A_H(x)|}{2^d}\right)^\alpha.
\]
由于 $0<\alpha<1$，函数 $t^\alpha$ 为凹函数。Jensen 不等式以及第一步的随机桶密度上界说明
\[
  \E_H\E_x
  \left(\frac{|A_H(x)|}{2^d}\right)^\alpha
  \le
  \left(
    \E_H\E_x\frac{|A_H(x)|}{2^d}
  \right)^\alpha
  \le
  \left(
    q+\exp\left(-\frac{(d/2-R)^2}{d}\right)
  \right)^\alpha.
\]
这就严格闭合了同一碰撞事件的下界与上界，并得到论文的核心命题。

\subsection{第五步：从 Hamming 立方体推广到 \texorpdfstring{$\ell_s$}{l-s} 空间}

由于 $\{0,1\}^d$ 可以看作 $\ell_s^d$ 的子集，并且在 0-1 向量上距离满足
\[
  \|x-y\|_s = \|x-y\|_1^{1/s},
\]
Hamming 距离上的下界可以转化为 $\ell_s$ 空间中的下界。最终得到：
\[
  \rho_s(c)
  \ge
  \frac{e^{1/c^s}-1}{e^{1/c^s}+1}
  \ge
  \frac{0.462}{c^s}.
\]

原因是，在 0-1 向量上，$\ell_s$ 距离等于 Hamming 距离的 $1/s$ 次方。如果在 Hamming 距离中远近阈值的比值是 $c^s$，
那么在 $\ell_s$ 距离中对应的比值就是 $c$。因此，Hamming 立方体上的下界可以转化为 $\ell_s$ 空间中的下界。
参数选择还需要保证 Hamming 球体积估计中的误差项消失，因此不能只写成 $R\approx d/2$，更不能直接取 $R=d/2$。
一个具体选择是：对于充分大的 $d$，取 $r_d$ 为不超过
\[
  \frac{1}{c^s}
  \left(\frac d2-\sqrt{d\log d}\right)
\]
的最大奇数，并令 $R_d=c^s r_d$。于是
\[
  R_d
  =
  \frac d2-\sqrt{d\log d}+O(1)
  =
  \frac d2-\Theta(\sqrt{d\log d}),
  \qquad
  r_d=\frac{R_d}{c^s}.
\]
这个选择同时保证 $R_d<d/2$、$r_d$ 为奇数，且
\[
  \exp\left(-\frac{(d/2-R_d)^2}{d}\right)
  =d^{-1+o(1)}
  \longrightarrow 0.
\]
另一方面，
\[
  \frac{2r_d}{d}\longrightarrow\frac{1}{c^s},
\]
所以随机游走引理中的指数满足
\[
  \frac{e^{2r_d/d}-1}{e^{2r_d/d}+1}
  \longrightarrow
  \frac{e^{1/c^s}-1}{e^{1/c^s}+1}.
\]
令 $d\to\infty$ 后，误差项消失，就得到论文的主定理。

\section{正确性与优越性}

\subsection{证明技术的正确性}

证明的正确性首先依赖于对随机性的严格区分。远点条件中的概率来自随机抽取
$H\sim\mathcal H$，因此它控制的是随机桶 $A_H(x)$ 大小的期望，而不是每个固定哈希函数的桶大小。
近点条件同时涉及随机哈希函数和随机游走终点，准确的表述是
\[
  \Prb_{H\sim\mathcal H,\,W_r}
  [H(W_r(x))=H(x)]\ge p.
\]
对每个可能的游走终点先应用 LSH 近点条件，再对 $W_r$ 取平均，即可得到该联合概率下界，
并未对固定的 $H$ 作额外假设。

其次，证明比较的是同一个碰撞事件。LSH 定义从下方以 $p$ 界定该事件；条件于哈希桶后，
随机游走引理和 Bonami--Beckner 不等式从上方界定它。由于
\[
  0<\alpha=\frac{e^{2r/d}-1}{e^{2r/d}+1}<1,
\]
$t^\alpha$ 的凹性保证 Jensen 不等式的方向正确，从而得到核心命题。最后，所选参数使 Hamming 球误差项趋于零，
而 $\{0,1\}^d\subseteq\ell_s^d$ 及
$\|x-y\|_s=\|x-y\|_1^{1/s}$ 保证子空间下界能够合法推广到整个 $\ell_s^d$。
因此，概率空间、上下界比较、极限过程和空间推广四个环节相互衔接，没有引入定义之外的假设。

\subsection{证明技术的优越性}

这里的“优越性”不是指论文提出了更快的近似最近邻算法，而是指其下界方法具有更强的普适性。
对某个具体 LSH 构造的分析只能说明该构造的性能；本文直接从敏感哈希族的定义出发，
所以结论同时约束所有数据无关的标准 LSH 哈希族。这把“尚未找到更好的构造”转化为
“该模型内不存在数量级更好的构造”。

在定量上，本文对 $\ell_1$ 和 $\ell_2$ 分别给出 $\Omega(1/c)$ 与 $\Omega(1/c^2)$ 下界，
同当时已知的 $O(1/c)$ 与 $O(1/c^2)$ 上界在渐近量级上匹配。更重要的是，随机桶大小与随机游走稳定性之间的矛盾
揭示了下界产生的结构性原因，而不只是给出一个数值不等式。这套 Hamming 立方体、Fourier 分析和噪声稳定性方法
也为后续 O'Donnell--Wu--Zhou 的紧下界提供了基础。不过，本文仍存在常数差距，其历史地位和后续改进将在下一章具体比较。

\section{与已有及后续工作的比较}

\subsection{已有上界与本文下界}

在本文之前，Indyk 和 Motwani 的构造在 $\ell_1$ 中达到
\[
  \rho_1(c)\le \frac{1}{c}.
\]
Datar 等人的 $p$-稳定分布方法进一步处理了 $\ell_s$ 空间 \cite{datar2004}，Andoni 和 Indyk
则在 $\ell_2$ 中得到 $\rho_2(c)\le 1/c^2$ \cite{andoni2005}。这些结果是构造上界；
Motwani、Naor 和 Panigrahy 首次给出可与之比较的统一下界。在 $\ell_1$ 中，本文证明
\[
  \rho_1(c)
  \ge
  \frac{e^{1/c}-1}{e^{1/c}+1}
  \sim
  \frac{1}{2c}
  \qquad (c\to\infty).
\]
这把差距缩小到常数因子，但并未得到标准数据无关 LSH 的最终紧界。

\subsection{O'Donnell--Wu--Zhou 紧下界}

O'Donnell、Wu 和 Zhou 后来利用布尔函数噪声稳定性的对数凸性，把 Hamming 空间中的下界提高为
\[
  \rho\ge \frac{1}{c}-o_d(1)
\]
\cite{odonnell2011}。在 $q\ge 2^{-o(d)}$（包括固定 $q$）的通常范围内，该结果与 $1/c$ 上界匹配，
并推出一般 $\ell_s$ 空间的 $1/c^s-o_d(1)$ 下界，特别是 Euclidean 空间中的紧 $1/c^2$ 下界。
因此，本文应评价为开创并奠定了标准 LSH 下界方法，而不是给出了该模型的最终紧常数。

\subsection{数据依赖哈希}

另一条路线是数据依赖哈希。Andoni 和 Razenshteyn 构造了 Hamming 空间中
\[
  \rho=\frac{1}{2c-1}+o(1)
\]
的方案，Euclidean 空间中的相应指数为 $1/(2c^2-1)$ \cite{andoni2015datadependent}。
他们又在限制哈希函数表示复杂度的简洁模型中证明匹配下界；例如 $\ell_1$ 中有
\[
  \rho\ge\frac{1}{2c-1}-o(1)
\]
\cite{andoni2016lower}。这说明改变模型可以严格优于数据无关 LSH，但新模型也有自己的紧边界。

\section{评价、局限与总结}

\subsection{贡献与方法价值}

本文最重要的贡献有两点。第一，它不再逐个分析具体构造，而是直接从“近点高碰撞、远点低碰撞”的定义推出
整个标准 LSH 模型的下界。第二，它把算法问题转化为 Hamming 立方体上的小集合扩散问题，
以随机桶、随机游走、Fourier 展开和 Bonami--Beckner 不等式形成了一条可复用的证明路线。
这种转化既给出定量结论，也解释了下界产生的原因。

\subsection{适用范围与局限}

论文也有明确局限。其结论只约束数据无关的标准 LSH，并不是所有近似最近邻数据结构的下界；
它对数据依赖哈希、图方法及其他计算模型不直接适用。此外，本文在 $\ell_1$ 中得到的渐近常数约为 $1/2$，
与已有 $1/c$ 上界仍有常数差距，后来才由 O'Donnell--Wu--Zhou 闭合。最后，主定理是高维渐近结论，
不能直接当作有限维实际系统的精确性能预测。

\subsection{总结}

总体而言，Motwani、Naor 和 Panigrahy 的论文首次证明高维 $\ell_s$ 空间中的标准 LSH 必须满足
\[
  \rho_s(c)
  \ge
  \frac{e^{1/c^s}-1}{e^{1/c^s}+1}
  =
  \Omega(1/c^s).
\]
它的持久价值不在于最终常数，而在于揭示了随机桶大小与随机游走稳定性之间的结构性矛盾，
并为后续紧下界奠定了方法基础。

\begin{thebibliography}{9}

\bibitem{mnp2006}
R. Motwani, A. Naor, and R. Panigrahy.
\newblock Lower Bounds on Locality Sensitive Hashing.
\newblock \emph{Proceedings of the Twenty-second Annual Symposium on Computational Geometry}, 2006.

\bibitem{im1998}
P. Indyk and R. Motwani.
\newblock Approximate Nearest Neighbors: Towards Removing the Curse of Dimensionality.
\newblock \emph{Proceedings of STOC}, 1998.

\bibitem{datar2004}
M. Datar, N. Immorlica, P. Indyk, and V. S. Mirrokni.
\newblock Locality-Sensitive Hashing Scheme Based on p-Stable Distributions.
\newblock \emph{Proceedings of the Twentieth Annual Symposium on Computational Geometry}, 2004.

\bibitem{andoni2005}
A. Andoni and P. Indyk.
\newblock Faster Algorithms for High-Dimensional Nearest Neighbor Search.
\newblock Manuscript, 2005.

\bibitem{odonnell2011}
R. O'Donnell, Y. Wu, and Y. Zhou.
\newblock Optimal Lower Bounds for Locality-Sensitive Hashing (Except When $q$ Is Tiny).
\newblock \emph{Proceedings of ITCS}, 2011.

\bibitem{andoni2015datadependent}
A. Andoni and I. Razenshteyn.
\newblock Optimal Data-Dependent Hashing for Approximate Near Neighbors.
\newblock \emph{Proceedings of STOC}, 2015.

\bibitem{andoni2016lower}
A. Andoni and I. Razenshteyn.
\newblock Tight Lower Bounds for Data-Dependent Locality-Sensitive Hashing.
\newblock \emph{Proceedings of the 32nd International Symposium on Computational Geometry}, 2016.

\end{thebibliography}

\end{document}
